% BEGIN PROFESSOR FEEDBACK PF30: Conclusion
\section{Conclusion}

Malicious activity on Ethereum is difficult to detect from individual
transactions. An address can look ordinary on its own even when it is part of a
larger sequence involving several addresses. This thesis examined whether
temporal motifs make such patterns useful for detection. Each address was
represented by the roles it occupies in ordered three-edge patterns that occur
within a temporal window $\delta$. This representation was then compared with
conventional transaction aggregates and static graph features.

The main contribution is an evaluation of temporal motif fingerprints
on labelled Ethereum addresses. The evaluation covers binary detection on the
HuggingFace Ethereum Fraud Dataset and a separate Real-CATS sample, as well as
category classification across nine Real-CATS behaviour labels. The pipeline
counts the retained motif roles for several temporal windows, keeps the five
Real-CATS transaction layers separate, and accounts for transfers that share a
timestamp. All main comparisons use held-out addresses. This makes it possible
to test whether a visible motif signal also improves prediction on addresses
that were not used to select the model or temporal window.
The available labels identify malicious addresses and reported behaviour
categories rather than independently verified coordinated campaigns. The
evaluation therefore measures address-level behavioural signal, not direct
campaign detection.

The answer to RQ1 is that temporal motifs contain information, but do not
reliably improve malicious-address detection beyond the conventional baseline.
Adding the fingerprint increases average precision in all six HuggingFace
cohort comparisons, although only three of these increases are statistically
significant. In the 40-role HuggingFace analysis, the increases are 0.0237 in
the controlled cohort and 0.0161 in the prevalence cohort, with both confidence
intervals including zero. The binary Real-CATS comparison produces an increase
of only 0.0011 and a lower F1 score. At the same time, the role-by-role
fingerprint performs better than motif volume alone. Its advantage in the
40-role analysis is 0.1173 average precision in the controlled cohort and
0.1528 in the prevalence cohort. The way an address participates in motifs is
therefore more informative than the total number of motifs, but much of this
information overlaps with simpler address features.

For RQ2, increasing the temporal window generally gives higher average precision
and provides a non-zero
fingerprint for more addresses. In the 40-role analysis, motif
coverage rises from 12.2\% to 77.7\% in the controlled cohort and from
11.1\% to 75.9\% in the prevalence cohort between five minutes and 30 days.
The 30-day window gives
the highest development score in both cohorts, but two weeks remains within one
standard error of the maximum. For no HuggingFace sample is the selected maximum
significantly better than every other tested window under all timestamp
orderings. Since 30
days is also the largest tested value, it cannot be treated as a universal
optimum, and no unique optimum within the tested grid is established either.
The selected window describes the maximum span of a counted candidate pattern,
not the duration of a malicious
campaign.

RQ3 shows that the nine Real-CATS behaviour categories have detectable
differences in their motif fingerprints. At the shared two-week window, the
classifier reaches a macro average precision of 0.1953, compared with an
average class share of 0.1111. The permutation result remains significant under
all seven timestamp orderings. However, the macro F1 score is only 0.1914 and
the categories overlap strongly in the confusion matrix. The fingerprints
therefore contain category information, but they are not sufficient for
accurate category classification on their own. RQ3 establishes separability
within this representation, not superiority over conventional behavioural
features for category classification.

The answer to RQ4 is no. Replacing the shared two-week window with separately
selected category windows changes held-out macro average precision by only
0.0009. The confidence interval includes zero, and none of the category-level
changes is significant after correction for multiple comparisons. Most of the
selected windows also change when equal-timestamp transfers are ordered
differently. These development selections should not be interpreted as stable
time scales of the corresponding behaviour categories.

The central result of the thesis is therefore a distinction between behavioural
information and additional predictive value. Temporal motif fingerprints reveal
how an address participates in local, ordered transaction patterns. This signal
is visible in motif-only detection and in the comparison between Real-CATS
categories. It does not consistently improve a detector that already has access
to strong transaction aggregates and static graph features. Temporal motifs are
therefore useful as a description of address behaviour, but the results do not
support using them as a replacement for conventional features or as a generally
reliable improvement to them.
% END PROFESSOR FEEDBACK PF30

